\section{Results}
\label{sec:results}

\subsection{Experimental campaign overview}
\label{subsec:results-overview}

The definitive campaign contained 1272 scheduled experimental tasks
across the four research questions. This section reports only measured
runs when evaluating the research hypotheses; warm-up executions are
reported separately from the principal statistical evidence.

The results are presented in the order of the research questions.
Incomplete executions are retained according to the censoring policy
defined in Section~\ref{subsec:censoring}.

\subsection{RQ1: Bounded property verification}
\label{subsec:rq1-results}

RQ1 asks whether the valid formal models satisfy the seven declared
protocol properties within the evaluated bounds.

Across both formal tools and the three bound profiles, RQ1 contained
420 measured property runs. Of these, 350 completed and 70 were
censored by timeout. No completed run produced a counterexample.

Table~\ref{tab:rq1-summary} summarizes the result by tool and bound
profile.

\begin{table}[htbp]
\centering
\caption{RQ1 bounded property-verification results by tool and profile.}
\label{tab:rq1-summary}
\begin{tabular}{llrrrr}
\hline
Tool & Profile & Measured & Completed & No counterexample & Timeout \\
\hline
Alloy & Small  & 70 & 70 & 70 & 0  \\
Alloy & Medium & 70 & 70 & 70 & 0  \\
Alloy & Large  & 70 & 70 & 70 & 0  \\
TLC   & Small  & 70 & 70 & 70 & 0  \\
TLC   & Medium & 70 & 70 & 70 & 0  \\
TLC   & Large  & 70 & 0  & 0  & 70 \\
\hline
Total &        & 420 & 350 & 350 & 70 \\
\hline
\end{tabular}
\end{table}

For Alloy, all seven properties completed without a counterexample
under the small, medium, and large profiles. Each property-profile
combination was measured ten times, yielding 210 completed Alloy runs.

For TLC, all seven properties completed without a counterexample under
the small and medium profiles, yielding 140 completed runs. In
contrast, every measured execution under the large TLC profile reached
the predefined 1800-second timeout. This produced 70 censored
observations, corresponding to seven properties and ten measured
repetitions per property.

No out-of-memory or tool-error outcome was observed in the RQ1
measurements. The incomplete evidence is therefore concentrated
entirely in the large TLC profile.

\subsubsection{Answer to RQ1}
\label{subsubsec:rq1-answer}

Within the 350 completed bounded model-checking runs, no violation of
the seven evaluated properties was observed. This provides supporting
evidence for H1 within the completed configurations.

However, the 70 large-profile TLC measurements did not complete within
the predefined resource limit. Their absence of a counterexample
cannot be interpreted as successful property verification.

H1 is therefore partially supported under censoring. The result is
restricted to the completed bounded configurations and does not
constitute an unbounded proof of protocol correctness.

\subsection{RQ2: Mutation-based validation}
\label{subsec:rq2-results}

RQ2 evaluates whether the declared property suite detects predefined
scientific mutations that remove controls from the cross-shard
protocol.

The experiment contained ten scientific mutants: five instantiated in
the TLA+ model and five in the Alloy model. Each mutant was evaluated
through ten measured repetitions at the medium profile, producing 100
measured mutation runs.

All ten scientific mutants were detected through their designated
target properties, and detection remained consistent across all
measured repetitions. The resulting mutation score was

\begin{equation}
\mathrm{MutationScore}
=
\frac{10}{10}
=
1.0.
\label{eq:rq2-mutation-score}
\end{equation}

Table~\ref{tab:rq2-summary} summarizes the aggregate result.

\begin{table}[htbp]
\centering
\small
\caption{RQ2 scientific mutants, target properties, and detection results.}
\label{tab:rq2-summary}
\begin{tabular}{llllr}
\hline
Tool & Mutant & Target property & Detected & Runs \\
\hline
Alloy &
commit-after-abort &
TerminalStateIrreversibility &
Yes & 10/10 \\

Alloy &
credit-before-receipt &
DestinationCreditRequiresValidReceipt &
Yes & 10/10 \\

Alloy &
no-replay &
NoReceiptReplay &
Yes & 10/10 \\

Alloy &
quorum-bypass &
QuorumRequired &
Yes & 10/10 \\

Alloy &
timeout-without-release &
EventuallyReleasedAfterTimeout &
Yes & 10/10 \\

TLC &
commit-after-abort &
DecisionConsistency &
Yes & 10/10 \\

TLC &
credit-before-receipt &
DestinationCreditRequiresValidReceipt &
Yes & 10/10 \\

TLC &
no-replay &
NoReceiptReplay &
Yes & 10/10 \\

TLC &
quorum-bypass &
QuorumRequired &
Yes & 10/10 \\

TLC &
timeout-without-release &
EventuallyReleasedAfterTimeout &
Yes & 10/10 \\
\hline
\end{tabular}
\end{table}

The commit-after-abort defect is associated with different designated
target properties in the two formal representations:
\texttt{TerminalStateIrreversibility} in Alloy and
\texttt{DecisionConsistency} in TLA+. This difference follows the
tool-specific encodings of terminal consistency and does not affect
the mutant-level detection criterion.

The evaluated defect classes covered removal of replay protection,
destination credit without the required receipt condition, commit after
an abort-like terminal decision, timeout without release of source
funds, and quorum bypass.

Detection was evaluated at the mutant level rather than by counting
individual counterexample-producing repetitions as separate mutants.
Therefore, the 100 measured runs represent repeated evaluations of ten
scientific mutants, not 100 independent mutant definitions.

The repeated measurements also verified consistency of the logical
outcome: each mutant continued to expose its designated target-property
violation across all ten measured repetitions.

No absolute execution-time comparison between TLC and Alloy is used to
support RQ2. The two formal representations have different exploration
semantics, and the relevant result for this research question is
whether each predefined mutant exposes its designated property
violation.

\subsubsection{Answer to RQ2}
\label{subsubsec:rq2-answer}

All ten predefined scientific mutants were detected through their
target properties, yielding a mutation score of 1.0 with consistent
detection across the measured repetitions.

H2 is therefore supported for the predefined scientific mutant
catalogue.

This result demonstrates that the evaluated property suite is
sensitive to the injected defect classes considered in the experiment.
It does not establish completeness with respect to arbitrary
cross-shard protocol defects or mutations outside the predefined
catalogue.
